\clearpage\SourceLeaf{520}
\renewcommand{\BookRunningTitle}{Achter Abschnitt.}
\pdfbookmark[0]{Achter Abschnitt. Ausführlichere Betrachtung einiger Untersuchungen zur Kreisteilung.}{sectio-octava}
\begin{center}
\textbf{ACHTER ABSCHNITT.}\\[1em]
AUSFÜHRLICHERE BETRACHTUNG EINIGER UNTERSUCHUNGEN ZUR KREISTEILUNG.
\end{center}
\BookRule{.18\textwidth}
\par\medskip\Needspace{4\baselineskip}\pdfbookmark[1]{367.}{sectio-octava-art-367}\centerline{\textbf{367.}}\nobreak\smallskip
Was wir im späteren Teil des siebten Abschnitts von Artikel $355$ an dargelegt haben, gibt gewiss gewichtige Beispiele für die große Fruchtbarkeit der Theorie der Kreisteilung und für den erstaunlichen Zusammenhang, der dieses Gebiet mit verschiedenen \textit{arithmetischen} Untersuchungen verbindet. Wegen der allzu engen Grenzen von Raum und Zeit konnten wir dieses Feld dort jedoch nur leicht berühren; je weiter wir darin vordringen, desto reicher ist die Ernte, mit der es unsere Bemühungen belohnt. Deshalb wollen wir hier die eine oder andere der dort begonnenen Fragen erneut aufnehmen und ausführlicher behandeln: Die Leser werden gewiss mit großem Staunen sehen, dass die Lösung mehrerer arithmetischer Probleme, die jeder für himmelweit von diesem Gebiet entfernt gehalten hätte, auf dieser Grundlage beruht.
\par\medskip\Needspace{4\baselineskip}\pdfbookmark[1]{368.}{sectio-octava-art-368}\centerline{\textbf{368.}}\nobreak\smallskip
Ein besonders fruchtbares Thema bietet die in Artikel $356$ begonnene Untersuchung. Dort teilten wir die Gesamtheit der Wurzeln der Gleichung $x^n-1=0$ (ohne eins) in zwei Klassen und zeigten, wie die Summe in jeder Klasse zu bestimmen ist. Dabei ergaben sich $=-\frac12+\frac12\surd n$ und $-\frac12-\frac12\surd n$ für den Fall, dass $n$ die Form $4n+1$ hat, oder $=-\frac12+\frac12\surd -n$ und $-\frac12-\frac12\surd -n$ für den Fall, dass $n$ die Form $4n+3$ hat. Wir hatten uns dort aber nicht nur \SourceLeaf{521}auf den Fall beschränkt, dass $n$ eine \textit{Primzahl} ist, sondern auch, was noch viel schwerer wog, das \textit{Vorzeichen} des Wurzelausdrucks unbestimmt gelassen oder vielmehr versäumt, diese mit wenigen Worten angedeutete Bestimmung durch einen festen Beweis zu stützen. Diese Mängel müssen wir daher vor allem anderen beheben.
\par\medskip\Needspace{4\baselineskip}\pdfbookmark[1]{369.}{sectio-octava-art-369}\centerline{\textbf{.369.}}\nobreak\smallskip
Nun sei $n$ eine beliebige positive ganze Zahl und $R$ eine Wurzel der Gleichung $x^n-1\pdot=0$, deren Potenzen vor der Potenz mit Exponenten $n$ sämtlich von eins verschieden sind (siehe Artikel $359$, II.). Wie in Abschnitt VII bezeichnen wir mit $[\lambda]$ die Potenz $R^\lambda$, sodass $[0]=1$, $[1]$, $[2]$, $[3]\dotsfour[n-1]$ alle Wurzeln der Gleichung $x^n-1=0$ darstellen. Ferner bezeichnen wir die Summe
\[
{}[0]+[1]+[4]+[9]+\dotsfour+[(n-1)^2]\quad\text{mit }\Sigma[\mathfrak{Q}]
\]
und allgemeiner
\[
{}[0]+[\lambda]+[4\lambda]+[9\lambda]\dotsfour+[\lambda(n-1)^2]\quad\text{mit }\Sigma[\mathfrak{Q}\lambda]
\]
Dabei bezeichnet $\mathfrak{Q}$ ohne nähere Festlegung die Quadrate der Zahlen $0,1,2,3\dotsfour n-1$. Wie allgemein $[\lambda]=[\mu]$ gilt, wenn $\lambda$, $\mu$ beliebige (positive oder negative) ganze Zahlen sind, die modulo $n$ kongruent sind, so gilt offenbar auch $\Sigma[\mathfrak{Q}\lambda]=\Sigma[\mathfrak{Q}\mu]$, wenn $\lambda\equiv\mu$. Nach diesen Vorbereitungen erhalten wir Folgendes.
\par\medskip\Needspace{4\baselineskip}\pdfbookmark[1]{370.}{sectio-octava-art-370}\centerline{\textbf{370.}}\nobreak\smallskip
\textsc{Aufgabe.} \textit{Das Produkt der beiden Summen $\Sigma[\mathfrak{Q}]$ und $\Sigma[-\mathfrak{Q}]$ bestimmen.}
\par\noindent\textit{Lösung.} Wegen $nn\equiv0$, $(n+1)^2\equiv1$, $(n+2)^2\equiv4$ usw. $(\mathrm{mod}.n)$ ergibt sich leicht, dass $\Sigma[\mathfrak{Q}]$ wird
\[
\begin{aligned}
&=[1]+[4]+[9]+[16]\dotsfive+[nn]\\
&=[4]+[9]+[16]+[25]\dotsfive+[(n+1)]^2\\
&=[9]+[16]+[25]+[36]\dotsfive+[(n+2)^2]
\end{aligned}
\]
\centerline{usw.\quad oder allgemein}
\[
{}=[kk]+[(k+1)^2]+[(k+2)^2]+[(k+3)^2]\dotsfour+[(n+k-1)^2]
\]
Daraus folgt für $[-kk]\times\Sigma[\mathfrak{Q}]$
\[
{}=[0]+[2k+1]+[4k+4]+[6k+9]\dotsfive+[(n-1)^2+2(n-1)k]
\]
